%%%%%%%%%%%%%%%%%%%%%part-01.tex%%%%%%%%%%%%%%%%%%%%%%%%%%%%%%%%%%
% 
% sample part title
%
% Use this file as a template for your own input.
%
%%%%%%%%%%%%%%%%%%%%%%%% Springer %%%%%%%%%%%%%%%%%%%%%%%%%%

\begin{partbacktext}
\part{Fundamentos axiomáticos y fenomenología de sistemas dinámicos discretos}

\noindent Esta primera parte del documento establece las bases teóricas y axiomáticas para el estudio de la computabilidad en autómatas celulares complejos. Se aborda la computación no convencional desde la perspectiva de los sistemas dinámicos discretos, definiendo las propiedades matemáticas que permiten emplear dinámicas naturales para el procesamiento de información. El objetivo principal de esta sección no es únicamente introducir conceptos, sino estructurar el problema desde una perspectiva formal antes de proceder a su implementación algorítmica.

En el Capítulo \ref{cap:preliminares}, se define formalmente el objeto de estudio y se plantea el problema computacional. Se establece el contexto fenomenológico en la frontera entre la complejidad topológica, la reversibilidad y la computabilidad universal. Asimismo, se delimitan las hipótesis sobre la emergencia en reglas de evolución específicas, como \textit{Spiral} y \textit{Beehive}, y se presentan los criterios de demostración analíticos y sintéticos (Sección \ref{sec:objetivos}) que regirán la validación del modelo propuesto.

Posteriormente, en el Capítulo \ref{cap:teoria_formal}, se construye la teoría matemática de los autómatas celulares. A partir de fundamentos de teoría de conjuntos y topología discreta, se axiomatizan los componentes fundamentales del sistema: la retícula espacial, las vecindades acotadas y las funciones de transición. Se analiza la dinámica global del espacio de configuraciones y se formaliza el papel de la irreversibilidad. Este último concepto es crucial, ya que la disipación de estados es un requerimiento matemático indispensable para el diseño de operaciones lógicas deterministas mediante colisiones.

En conjunto, los axiomas y definiciones presentados en esta parte establecen el terreno lógico estrictamente necesario para el desarrollo del modelo. Esta fundamentación es la base sobre la cual se construirá la representación analítica del Capítulo \ref{cap:geometria_computacional}, la demostración de transformaciones y mapeos del Capítulo \ref{chap:isomorfismos}, y el análisis cinemático y de complejidad algorítmica del Capítulo \ref{cap:dinamica_particulas}.

\end{partbacktext}